\subsection*{Вариант 1 для  Рыльской Эвилины Ильгамовны. Дроби \\ \today   \currenttime} 
  \begin{enumerate}
\item Найди число, если его $\frac{1}{5}$  часть равна 25. Нарисуй схему.  \\ \nopagebreak \begin{tikzpicture}
 \draw[step=0.7cm ,gray,very thin] (0,0) grid (18,6);
 % Обводка внешнего контура  \draw[black, thick] (0,0) rectangle (18,6);
 \end{tikzpicture}
\item Найди $\frac{1}{8}$ часть от 88. Нарисуй схему. \\ \nopagebreak \begin{tikzpicture}
 \draw[step=0.7cm ,gray,very thin] (0,0) grid (18,6);
 % Обводка внешнего контура  \draw[black, thick] (0,0) rectangle (18,6);
 \end{tikzpicture}
\item Найди число, если его $\frac{5}{11}$  часть равна 55. Нарисуй схему.  \\ \nopagebreak \begin{tikzpicture}
 \draw[step=0.7cm ,gray,very thin] (0,0) grid (18,6);
 % Обводка внешнего контура  \draw[black, thick] (0,0) rectangle (18,6);
 \end{tikzpicture}
\item Найди $\frac{7}{12}$ часть от 84.  Нарисуй схему. \\ \nopagebreak \begin{tikzpicture}
 \draw[step=0.7cm ,gray,very thin] (0,0) grid (18,6);
 % Обводка внешнего контура  \draw[black, thick] (0,0) rectangle (18,6);
 \end{tikzpicture}
\item Найди число, если его  3 \%  равны 1800. Нарисуй схему.  \\ \nopagebreak \begin{tikzpicture}
 \draw[step=0.7cm ,gray,very thin] (0,0) grid (18,6);
 % Обводка внешнего контура  \draw[black, thick] (0,0) rectangle (18,6);
 \end{tikzpicture}
\item Найди 12 \%  от 15600.  Нарисуй схему. \\ \nopagebreak \begin{tikzpicture}
 \draw[step=0.7cm ,gray,very thin] (0,0) grid (18,6);
 % Обводка внешнего контура  \draw[black, thick] (0,0) rectangle (18,6);
 \end{tikzpicture}
\end{enumerate}  
\newpage
\subsection*{Вариант 2 для  Абелюк Марии. Дроби \\ \today   \currenttime} 
  \begin{enumerate}
\item Найди число, если его $\frac{1}{10}$  часть равна 60. Нарисуй схему.  \\ \nopagebreak \begin{tikzpicture}
 \draw[step=0.7cm ,gray,very thin] (0,0) grid (18,6);
 % Обводка внешнего контура  \draw[black, thick] (0,0) rectangle (18,6);
 \end{tikzpicture}
\item Найди $\frac{1}{7}$ часть от 28. Нарисуй схему. \\ \nopagebreak \begin{tikzpicture}
 \draw[step=0.7cm ,gray,very thin] (0,0) grid (18,6);
 % Обводка внешнего контура  \draw[black, thick] (0,0) rectangle (18,6);
 \end{tikzpicture}
\item Найди число, если его $\frac{6}{12}$  часть равна 72. Нарисуй схему.  \\ \nopagebreak \begin{tikzpicture}
 \draw[step=0.7cm ,gray,very thin] (0,0) grid (18,6);
 % Обводка внешнего контура  \draw[black, thick] (0,0) rectangle (18,6);
 \end{tikzpicture}
\item Найди $\frac{8}{14}$ часть от 112.  Нарисуй схему. \\ \nopagebreak \begin{tikzpicture}
 \draw[step=0.7cm ,gray,very thin] (0,0) grid (18,6);
 % Обводка внешнего контура  \draw[black, thick] (0,0) rectangle (18,6);
 \end{tikzpicture}
\item Найди число, если его  2 \%  равны 3200. Нарисуй схему.  \\ \nopagebreak \begin{tikzpicture}
 \draw[step=0.7cm ,gray,very thin] (0,0) grid (18,6);
 % Обводка внешнего контура  \draw[black, thick] (0,0) rectangle (18,6);
 \end{tikzpicture}
\item Найди 5 \%  от 9500.  Нарисуй схему. \\ \nopagebreak \begin{tikzpicture}
 \draw[step=0.7cm ,gray,very thin] (0,0) grid (18,6);
 % Обводка внешнего контура  \draw[black, thick] (0,0) rectangle (18,6);
 \end{tikzpicture}
\end{enumerate}  
\newpage
\subsection*{Вариант 3 для  Гондарева Ильи. Дроби \\ \today   \currenttime} 
  \begin{enumerate}
\item Найди число, если его $\frac{1}{7}$  часть равна 49. Нарисуй схему.  \\ \nopagebreak \begin{tikzpicture}
 \draw[step=0.7cm ,gray,very thin] (0,0) grid (18,6);
 % Обводка внешнего контура  \draw[black, thick] (0,0) rectangle (18,6);
 \end{tikzpicture}
\item Найди $\frac{1}{2}$ часть от 20. Нарисуй схему. \\ \nopagebreak \begin{tikzpicture}
 \draw[step=0.7cm ,gray,very thin] (0,0) grid (18,6);
 % Обводка внешнего контура  \draw[black, thick] (0,0) rectangle (18,6);
 \end{tikzpicture}
\item Найди число, если его $\frac{7}{14}$  часть равна 98. Нарисуй схему.  \\ \nopagebreak \begin{tikzpicture}
 \draw[step=0.7cm ,gray,very thin] (0,0) grid (18,6);
 % Обводка внешнего контура  \draw[black, thick] (0,0) rectangle (18,6);
 \end{tikzpicture}
\item Найди $\frac{12}{19}$ часть от 228.  Нарисуй схему. \\ \nopagebreak \begin{tikzpicture}
 \draw[step=0.7cm ,gray,very thin] (0,0) grid (18,6);
 % Обводка внешнего контура  \draw[black, thick] (0,0) rectangle (18,6);
 \end{tikzpicture}
\item Найди число, если его  15 \%  равны 9000. Нарисуй схему.  \\ \nopagebreak \begin{tikzpicture}
 \draw[step=0.7cm ,gray,very thin] (0,0) grid (18,6);
 % Обводка внешнего контура  \draw[black, thick] (0,0) rectangle (18,6);
 \end{tikzpicture}
\item Найди 6 \%  от 9000.  Нарисуй схему. \\ \nopagebreak \begin{tikzpicture}
 \draw[step=0.7cm ,gray,very thin] (0,0) grid (18,6);
 % Обводка внешнего контура  \draw[black, thick] (0,0) rectangle (18,6);
 \end{tikzpicture}
\end{enumerate}  
\newpage
\subsection*{Вариант 4 для  Аношина Матвея Антоновича. Дроби \\ \today   \currenttime} 
  \begin{enumerate}
\item Найди число, если его $\frac{1}{4}$  часть равна 20. Нарисуй схему.  \\ \nopagebreak \begin{tikzpicture}
 \draw[step=0.7cm ,gray,very thin] (0,0) grid (18,6);
 % Обводка внешнего контура  \draw[black, thick] (0,0) rectangle (18,6);
 \end{tikzpicture}
\item Найди $\frac{1}{13}$ часть от 182. Нарисуй схему. \\ \nopagebreak \begin{tikzpicture}
 \draw[step=0.7cm ,gray,very thin] (0,0) grid (18,6);
 % Обводка внешнего контура  \draw[black, thick] (0,0) rectangle (18,6);
 \end{tikzpicture}
\item Найди число, если его $\frac{8}{14}$  часть равна 112. Нарисуй схему.  \\ \nopagebreak \begin{tikzpicture}
 \draw[step=0.7cm ,gray,very thin] (0,0) grid (18,6);
 % Обводка внешнего контура  \draw[black, thick] (0,0) rectangle (18,6);
 \end{tikzpicture}
\item Найди $\frac{11}{18}$ часть от 198.  Нарисуй схему. \\ \nopagebreak \begin{tikzpicture}
 \draw[step=0.7cm ,gray,very thin] (0,0) grid (18,6);
 % Обводка внешнего контура  \draw[black, thick] (0,0) rectangle (18,6);
 \end{tikzpicture}
\item Найди число, если его  17 \%  равны 20400. Нарисуй схему.  \\ \nopagebreak \begin{tikzpicture}
 \draw[step=0.7cm ,gray,very thin] (0,0) grid (18,6);
 % Обводка внешнего контура  \draw[black, thick] (0,0) rectangle (18,6);
 \end{tikzpicture}
\item Найди 9 \%  от 5400.  Нарисуй схему. \\ \nopagebreak \begin{tikzpicture}
 \draw[step=0.7cm ,gray,very thin] (0,0) grid (18,6);
 % Обводка внешнего контура  \draw[black, thick] (0,0) rectangle (18,6);
 \end{tikzpicture}
\end{enumerate}  
\newpage
\subsection*{Вариант 5 для  Егоровой Мирославы Даниловны. Дроби \\ \today   \currenttime} 
  \begin{enumerate}
\item Найди число, если его $\frac{1}{7}$  часть равна 63. Нарисуй схему.  \\ \nopagebreak \begin{tikzpicture}
 \draw[step=0.7cm ,gray,very thin] (0,0) grid (18,6);
 % Обводка внешнего контура  \draw[black, thick] (0,0) rectangle (18,6);
 \end{tikzpicture}
\item Найди $\frac{1}{5}$ часть от 40. Нарисуй схему. \\ \nopagebreak \begin{tikzpicture}
 \draw[step=0.7cm ,gray,very thin] (0,0) grid (18,6);
 % Обводка внешнего контура  \draw[black, thick] (0,0) rectangle (18,6);
 \end{tikzpicture}
\item Найди число, если его $\frac{5}{12}$  часть равна 60. Нарисуй схему.  \\ \nopagebreak \begin{tikzpicture}
 \draw[step=0.7cm ,gray,very thin] (0,0) grid (18,6);
 % Обводка внешнего контура  \draw[black, thick] (0,0) rectangle (18,6);
 \end{tikzpicture}
\item Найди $\frac{5}{14}$ часть от 70.  Нарисуй схему. \\ \nopagebreak \begin{tikzpicture}
 \draw[step=0.7cm ,gray,very thin] (0,0) grid (18,6);
 % Обводка внешнего контура  \draw[black, thick] (0,0) rectangle (18,6);
 \end{tikzpicture}
\item Найди число, если его  16 \%  равны 20800. Нарисуй схему.  \\ \nopagebreak \begin{tikzpicture}
 \draw[step=0.7cm ,gray,very thin] (0,0) grid (18,6);
 % Обводка внешнего контура  \draw[black, thick] (0,0) rectangle (18,6);
 \end{tikzpicture}
\item Найди 11 \%  от 18700.  Нарисуй схему. \\ \nopagebreak \begin{tikzpicture}
 \draw[step=0.7cm ,gray,very thin] (0,0) grid (18,6);
 % Обводка внешнего контура  \draw[black, thick] (0,0) rectangle (18,6);
 \end{tikzpicture}
\end{enumerate}  
\newpage
\subsection*{Вариант 6 для  Карнауха Тимофея Павловича. Дроби \\ \today   \currenttime} 
  \begin{enumerate}
\item Найди число, если его $\frac{1}{11}$  часть равна 55. Нарисуй схему.  \\ \nopagebreak \begin{tikzpicture}
 \draw[step=0.7cm ,gray,very thin] (0,0) grid (18,6);
 % Обводка внешнего контура  \draw[black, thick] (0,0) rectangle (18,6);
 \end{tikzpicture}
\item Найди $\frac{1}{2}$ часть от 26. Нарисуй схему. \\ \nopagebreak \begin{tikzpicture}
 \draw[step=0.7cm ,gray,very thin] (0,0) grid (18,6);
 % Обводка внешнего контура  \draw[black, thick] (0,0) rectangle (18,6);
 \end{tikzpicture}
\item Найди число, если его $\frac{10}{19}$  часть равна 190. Нарисуй схему.  \\ \nopagebreak \begin{tikzpicture}
 \draw[step=0.7cm ,gray,very thin] (0,0) grid (18,6);
 % Обводка внешнего контура  \draw[black, thick] (0,0) rectangle (18,6);
 \end{tikzpicture}
\item Найди $\frac{3}{13}$ часть от 39.  Нарисуй схему. \\ \nopagebreak \begin{tikzpicture}
 \draw[step=0.7cm ,gray,very thin] (0,0) grid (18,6);
 % Обводка внешнего контура  \draw[black, thick] (0,0) rectangle (18,6);
 \end{tikzpicture}
\item Найди число, если его  8 \%  равны 6400. Нарисуй схему.  \\ \nopagebreak \begin{tikzpicture}
 \draw[step=0.7cm ,gray,very thin] (0,0) grid (18,6);
 % Обводка внешнего контура  \draw[black, thick] (0,0) rectangle (18,6);
 \end{tikzpicture}
\item Найди 4 \%  от 3200.  Нарисуй схему. \\ \nopagebreak \begin{tikzpicture}
 \draw[step=0.7cm ,gray,very thin] (0,0) grid (18,6);
 % Обводка внешнего контура  \draw[black, thick] (0,0) rectangle (18,6);
 \end{tikzpicture}
\end{enumerate}  
\newpage
\subsection*{Вариант 7 для  Костенко Каролины Романовны. Дроби \\ \today   \currenttime} 
  \begin{enumerate}
\item Найди число, если его $\frac{1}{8}$  часть равна 24. Нарисуй схему.  \\ \nopagebreak \begin{tikzpicture}
 \draw[step=0.7cm ,gray,very thin] (0,0) grid (18,6);
 % Обводка внешнего контура  \draw[black, thick] (0,0) rectangle (18,6);
 \end{tikzpicture}
\item Найди $\frac{1}{4}$ часть от 20. Нарисуй схему. \\ \nopagebreak \begin{tikzpicture}
 \draw[step=0.7cm ,gray,very thin] (0,0) grid (18,6);
 % Обводка внешнего контура  \draw[black, thick] (0,0) rectangle (18,6);
 \end{tikzpicture}
\item Найди число, если его $\frac{8}{13}$  часть равна 104. Нарисуй схему.  \\ \nopagebreak \begin{tikzpicture}
 \draw[step=0.7cm ,gray,very thin] (0,0) grid (18,6);
 % Обводка внешнего контура  \draw[black, thick] (0,0) rectangle (18,6);
 \end{tikzpicture}
\item Найди $\frac{7}{17}$ часть от 119.  Нарисуй схему. \\ \nopagebreak \begin{tikzpicture}
 \draw[step=0.7cm ,gray,very thin] (0,0) grid (18,6);
 % Обводка внешнего контура  \draw[black, thick] (0,0) rectangle (18,6);
 \end{tikzpicture}
\item Найди число, если его  18 \%  равны 10800. Нарисуй схему.  \\ \nopagebreak \begin{tikzpicture}
 \draw[step=0.7cm ,gray,very thin] (0,0) grid (18,6);
 % Обводка внешнего контура  \draw[black, thick] (0,0) rectangle (18,6);
 \end{tikzpicture}
\item Найди 10 \%  от 13000.  Нарисуй схему. \\ \nopagebreak \begin{tikzpicture}
 \draw[step=0.7cm ,gray,very thin] (0,0) grid (18,6);
 % Обводка внешнего контура  \draw[black, thick] (0,0) rectangle (18,6);
 \end{tikzpicture}
\end{enumerate}  
\newpage
\subsection*{Вариант 8 для  Минца Марка Евгеньевича. Дроби \\ \today   \currenttime} 
  \begin{enumerate}
\item Найди число, если его $\frac{1}{11}$  часть равна 88. Нарисуй схему.  \\ \nopagebreak \begin{tikzpicture}
 \draw[step=0.7cm ,gray,very thin] (0,0) grid (18,6);
 % Обводка внешнего контура  \draw[black, thick] (0,0) rectangle (18,6);
 \end{tikzpicture}
\item Найди $\frac{1}{7}$ часть от 49. Нарисуй схему. \\ \nopagebreak \begin{tikzpicture}
 \draw[step=0.7cm ,gray,very thin] (0,0) grid (18,6);
 % Обводка внешнего контура  \draw[black, thick] (0,0) rectangle (18,6);
 \end{tikzpicture}
\item Найди число, если его $\frac{9}{20}$  часть равна 180. Нарисуй схему.  \\ \nopagebreak \begin{tikzpicture}
 \draw[step=0.7cm ,gray,very thin] (0,0) grid (18,6);
 % Обводка внешнего контура  \draw[black, thick] (0,0) rectangle (18,6);
 \end{tikzpicture}
\item Найди $\frac{10}{20}$ часть от 200.  Нарисуй схему. \\ \nopagebreak \begin{tikzpicture}
 \draw[step=0.7cm ,gray,very thin] (0,0) grid (18,6);
 % Обводка внешнего контура  \draw[black, thick] (0,0) rectangle (18,6);
 \end{tikzpicture}
\item Найди число, если его  11 \%  равны 24200. Нарисуй схему.  \\ \nopagebreak \begin{tikzpicture}
 \draw[step=0.7cm ,gray,very thin] (0,0) grid (18,6);
 % Обводка внешнего контура  \draw[black, thick] (0,0) rectangle (18,6);
 \end{tikzpicture}
\item Найди 5 \%  от 7500.  Нарисуй схему. \\ \nopagebreak \begin{tikzpicture}
 \draw[step=0.7cm ,gray,very thin] (0,0) grid (18,6);
 % Обводка внешнего контура  \draw[black, thick] (0,0) rectangle (18,6);
 \end{tikzpicture}
\end{enumerate}  
\newpage
\subsection*{Вариант 9 для  Полякова Никиты. Дроби \\ \today   \currenttime} 
  \begin{enumerate}
\item Найди число, если его $\frac{1}{5}$  часть равна 15. Нарисуй схему.  \\ \nopagebreak \begin{tikzpicture}
 \draw[step=0.7cm ,gray,very thin] (0,0) grid (18,6);
 % Обводка внешнего контура  \draw[black, thick] (0,0) rectangle (18,6);
 \end{tikzpicture}
\item Найди $\frac{1}{5}$ часть от 40. Нарисуй схему. \\ \nopagebreak \begin{tikzpicture}
 \draw[step=0.7cm ,gray,very thin] (0,0) grid (18,6);
 % Обводка внешнего контура  \draw[black, thick] (0,0) rectangle (18,6);
 \end{tikzpicture}
\item Найди число, если его $\frac{6}{13}$  часть равна 78. Нарисуй схему.  \\ \nopagebreak \begin{tikzpicture}
 \draw[step=0.7cm ,gray,very thin] (0,0) grid (18,6);
 % Обводка внешнего контура  \draw[black, thick] (0,0) rectangle (18,6);
 \end{tikzpicture}
\item Найди $\frac{12}{18}$ часть от 216.  Нарисуй схему. \\ \nopagebreak \begin{tikzpicture}
 \draw[step=0.7cm ,gray,very thin] (0,0) grid (18,6);
 % Обводка внешнего контура  \draw[black, thick] (0,0) rectangle (18,6);
 \end{tikzpicture}
\item Найди число, если его  5 \%  равны 2000. Нарисуй схему.  \\ \nopagebreak \begin{tikzpicture}
 \draw[step=0.7cm ,gray,very thin] (0,0) grid (18,6);
 % Обводка внешнего контура  \draw[black, thick] (0,0) rectangle (18,6);
 \end{tikzpicture}
\item Найди 3 \%  от 6900.  Нарисуй схему. \\ \nopagebreak \begin{tikzpicture}
 \draw[step=0.7cm ,gray,very thin] (0,0) grid (18,6);
 % Обводка внешнего контура  \draw[black, thick] (0,0) rectangle (18,6);
 \end{tikzpicture}
\end{enumerate}  
\newpage
\subsection*{Вариант 10 для  Саргсян Эмилии Арамовны. Дроби \\ \today   \currenttime} 
  \begin{enumerate}
\item Найди число, если его $\frac{1}{5}$  часть равна 25. Нарисуй схему.  \\ \nopagebreak \begin{tikzpicture}
 \draw[step=0.7cm ,gray,very thin] (0,0) grid (18,6);
 % Обводка внешнего контура  \draw[black, thick] (0,0) rectangle (18,6);
 \end{tikzpicture}
\item Найди $\frac{1}{14}$ часть от 140. Нарисуй схему. \\ \nopagebreak \begin{tikzpicture}
 \draw[step=0.7cm ,gray,very thin] (0,0) grid (18,6);
 % Обводка внешнего контура  \draw[black, thick] (0,0) rectangle (18,6);
 \end{tikzpicture}
\item Найди число, если его $\frac{3}{11}$  часть равна 33. Нарисуй схему.  \\ \nopagebreak \begin{tikzpicture}
 \draw[step=0.7cm ,gray,very thin] (0,0) grid (18,6);
 % Обводка внешнего контура  \draw[black, thick] (0,0) rectangle (18,6);
 \end{tikzpicture}
\item Найди $\frac{12}{17}$ часть от 204.  Нарисуй схему. \\ \nopagebreak \begin{tikzpicture}
 \draw[step=0.7cm ,gray,very thin] (0,0) grid (18,6);
 % Обводка внешнего контура  \draw[black, thick] (0,0) rectangle (18,6);
 \end{tikzpicture}
\item Найди число, если его  3 \%  равны 3000. Нарисуй схему.  \\ \nopagebreak \begin{tikzpicture}
 \draw[step=0.7cm ,gray,very thin] (0,0) grid (18,6);
 % Обводка внешнего контура  \draw[black, thick] (0,0) rectangle (18,6);
 \end{tikzpicture}
\item Найди 10 \%  от 17000.  Нарисуй схему. \\ \nopagebreak \begin{tikzpicture}
 \draw[step=0.7cm ,gray,very thin] (0,0) grid (18,6);
 % Обводка внешнего контура  \draw[black, thick] (0,0) rectangle (18,6);
 \end{tikzpicture}
\end{enumerate}  
\newpage
\subsection*{Вариант 11 для  Строкатовой Киры Павловны. Дроби \\ \today   \currenttime} 
  \begin{enumerate}
\item Найди число, если его $\frac{1}{10}$  часть равна 40. Нарисуй схему.  \\ \nopagebreak \begin{tikzpicture}
 \draw[step=0.7cm ,gray,very thin] (0,0) grid (18,6);
 % Обводка внешнего контура  \draw[black, thick] (0,0) rectangle (18,6);
 \end{tikzpicture}
\item Найди $\frac{1}{8}$ часть от 64. Нарисуй схему. \\ \nopagebreak \begin{tikzpicture}
 \draw[step=0.7cm ,gray,very thin] (0,0) grid (18,6);
 % Обводка внешнего контура  \draw[black, thick] (0,0) rectangle (18,6);
 \end{tikzpicture}
\item Найди число, если его $\frac{7}{17}$  часть равна 119. Нарисуй схему.  \\ \nopagebreak \begin{tikzpicture}
 \draw[step=0.7cm ,gray,very thin] (0,0) grid (18,6);
 % Обводка внешнего контура  \draw[black, thick] (0,0) rectangle (18,6);
 \end{tikzpicture}
\item Найди $\frac{7}{13}$ часть от 91.  Нарисуй схему. \\ \nopagebreak \begin{tikzpicture}
 \draw[step=0.7cm ,gray,very thin] (0,0) grid (18,6);
 % Обводка внешнего контура  \draw[black, thick] (0,0) rectangle (18,6);
 \end{tikzpicture}
\item Найди число, если его  20 \%  равны 14000. Нарисуй схему.  \\ \nopagebreak \begin{tikzpicture}
 \draw[step=0.7cm ,gray,very thin] (0,0) grid (18,6);
 % Обводка внешнего контура  \draw[black, thick] (0,0) rectangle (18,6);
 \end{tikzpicture}
\item Найди 10 \%  от 12000.  Нарисуй схему. \\ \nopagebreak \begin{tikzpicture}
 \draw[step=0.7cm ,gray,very thin] (0,0) grid (18,6);
 % Обводка внешнего контура  \draw[black, thick] (0,0) rectangle (18,6);
 \end{tikzpicture}
\end{enumerate}  
\newpage
\subsection*{Вариант 12 для  Тавриной Марии Ивановны. Дроби \\ \today   \currenttime} 
  \begin{enumerate}
\item Найди число, если его $\frac{1}{9}$  часть равна 27. Нарисуй схему.  \\ \nopagebreak \begin{tikzpicture}
 \draw[step=0.7cm ,gray,very thin] (0,0) grid (18,6);
 % Обводка внешнего контура  \draw[black, thick] (0,0) rectangle (18,6);
 \end{tikzpicture}
\item Найди $\frac{1}{6}$ часть от 30. Нарисуй схему. \\ \nopagebreak \begin{tikzpicture}
 \draw[step=0.7cm ,gray,very thin] (0,0) grid (18,6);
 % Обводка внешнего контура  \draw[black, thick] (0,0) rectangle (18,6);
 \end{tikzpicture}
\item Найди число, если его $\frac{6}{14}$  часть равна 84. Нарисуй схему.  \\ \nopagebreak \begin{tikzpicture}
 \draw[step=0.7cm ,gray,very thin] (0,0) grid (18,6);
 % Обводка внешнего контура  \draw[black, thick] (0,0) rectangle (18,6);
 \end{tikzpicture}
\item Найди $\frac{12}{22}$ часть от 264.  Нарисуй схему. \\ \nopagebreak \begin{tikzpicture}
 \draw[step=0.7cm ,gray,very thin] (0,0) grid (18,6);
 % Обводка внешнего контура  \draw[black, thick] (0,0) rectangle (18,6);
 \end{tikzpicture}
\item Найди число, если его  16 \%  равны 32000. Нарисуй схему.  \\ \nopagebreak \begin{tikzpicture}
 \draw[step=0.7cm ,gray,very thin] (0,0) grid (18,6);
 % Обводка внешнего контура  \draw[black, thick] (0,0) rectangle (18,6);
 \end{tikzpicture}
\item Найди 11 \%  от 13200.  Нарисуй схему. \\ \nopagebreak \begin{tikzpicture}
 \draw[step=0.7cm ,gray,very thin] (0,0) grid (18,6);
 % Обводка внешнего контура  \draw[black, thick] (0,0) rectangle (18,6);
 \end{tikzpicture}
\end{enumerate}  
\newpage
